\section{Physical consistency}
\label{sec:consistency}

Five checks constrain the implementation against the physics it is built to
reproduce. Each check names its measured value.

\paragraph{Li\'enard--Wiechert kernel.}
The kernel of Section~\ref{sec:fielddef} is the present-time
Li\'enard--Wiechert field of a uniformly moving charge~\citep{jackson2007}.
The mesh reproduces the TSC-convolved Coulomb field of a single stationary
particle to a worst relative error of $4.88\times10^{-14}$,
E-component L2 residual $2.29\times10^{-25}$, and elementwise maximum
$2.28\times10^{-25}$, at $d=8,16,24,32$ along all six axis
directions, against the direct TSC stencil sum of the analytic Coulomb
kernel (Coulomb anchor, Section~\ref{sec:scaling}).

\paragraph{Lorentz contraction.}
The field of a uniformly moving charge is compressed into a pancake of
width $d/\gamma$ perpendicular to the motion. The seven-velocity
measurement (Appendix~\ref{app:sweep}, Fig.~\ref{fig:gp10}) tracks the
pancake: at $\gamma=10$ the mesh resolves the contracted field to
$1.2\times10^{-1}$ worst relative error on the mid-$z$ plane of a
$128^3$ grid with $5\times10^{-8}$\,m cells, and the error grows with
$\gamma$ because the pancake becomes thinner than the mesh spacing
($1.43$ and $3.21$ at $\gamma=50$ and $100$), not because of a solver
error. Raising the resolution from $64^3$ to $128^3$ at the same physical
domain leaves the relative errors unchanged.

\paragraph{Gauss's law.}
The discrete Gauss-law indicator, the Yee backward-difference
$\mathrm{div}\,E$ residual normalised to the peak deposited
$\rho/\varepsilon_0$, is $0.3534$ for the mesh alone on the $N=1000$ slice
geometry ($\gamma=1$) and $5.026$ for $\gamma=10$, against reference values
$0.0904$ and $0.0814$. With the near shell enabled it moves to $0.0895$ for
$\gamma=1$, reaching the reference value, and to $0.399$ for $\gamma=10$,
a $12.6\times$ drop that still leaves it $4.9\times$ above the reference
because the shell radius covers a fraction of the TSC support
(Appendix~\ref{app:slices}).

\paragraph{Yee staggering.}
The six field components are evaluated at the staggered positions of
Section~\ref{sec:readout}; the kernel sampling uses the same offsets, so the
convolution is already on the Yee grid and the read-out applies no
interpolation. The cross-validation against the PIConGPU
\texttt{referenceOnly} field uses identical Yee offsets, grid, and origin,
and the reference staggered point maps to a picBunch node by index with zero
snap deviation (Section~\ref{sec:validation}).

\paragraph{The $B = \boldsymbol{v}\times\boldsymbol{E}/c^2$ limit.}
The kernel satisfies $\mathbf{B}_n = \boldsymbol{\beta}\times\mathbf{E}_n$
exactly, and the mesh inherits it component by component. For a stationary
bunch ($\gamma=1$) every measured $B$ component is exactly zero
(Tables~\ref{tab:wholefield} and \ref{tab:multivel}); for motion along
$+y$ the component $B_y$ is identically zero in the measured fields
($B_y\equiv 0$ in Figs.~\ref{fig:valforwardB} and
\ref{fig:gp10}). Accordingly, the fully masked error panels in
Figs.~\ref{fig:gp1}--\ref{fig:gp100} (the whole $B$ row at $\gamma=1$;
the $B_y$ panel for $\gamma>1$), Fig.~\ref{fig:valforwardB} ($B_y$) and
Fig.~\ref{fig:largeB} ($B_x$) are fully masked, because the relative error
of an identically zero reference component is undefined. The magnetic field at
$\gamma=10$ is about
$3.3\!\times\!10^{-9}$ of the $E$-dominated local field
(Section~\ref{sec:accuracy}), which is why the
floored error metric is electric-only there (Section~\ref{sec:accuracy}).
